% 06_agents.tex — Multi-Agent Orchestration
% ============================================================

\section{Multi-Agent Orchestration}
\label{sec:agents}

The \pwm{} automates the full imaging pipeline --- from user intent to
reproducible result --- through a multi-agent architecture comprising
\textbf{9~agents} and \textbf{8~support classes}, totaling
\textbf{10,545~lines of Python}.  A critical design principle: \emph{all
agents run deterministically without requiring a large language model}.
An LLM may optionally be attached to generate narrative explanations,
but every gate decision, score, and recommendation is computed
mechanically from the YAML registries and Pydantic contracts.

% ============================================================
%  Pipeline overview
% ============================================================
\subsection{Pipeline Overview}

The end-to-end pipeline flow is as follows:
%
\begin{center}
\small
\begin{tabular}{@{}r@{\;$\rightarrow$\;}l@{}}
User Prompt       & \textbf{PlanAgent} \\
PlanAgent         & \textbf{PhotonAgent} $+$ \textbf{MismatchAgent}
                    $+$ \textbf{RecoverabilityAgent} \\
Gate Agents       & \textbf{AnalysisAgent} \\
AnalysisAgent     & \textbf{Negotiator} \\
Negotiator        & \textbf{PreFlightReportBuilder} \\
PreFlightReport   & \textbf{Pipeline Runner} \\
Pipeline Runner   & \textbf{RunBundle} \\
\end{tabular}
\end{center}

% ============================================================
%  Agent descriptions
% ============================================================
\subsection{Agent Descriptions}

\paragraph{PlanAgent.}
Parses user intent, maps the request to a registered modality, and
builds the \texttt{ImagingSystem} object.  All modality look-ups are
resolved against \texttt{modalities.yaml} --- no free-form strings are
accepted.

\paragraph{PhotonAgent.}
Computes the signal-to-noise ratio from the carrier budget (photon
count, detector quantum efficiency, read noise).  Classifies the noise
regime (shot-limited, read-limited, background-limited) and returns a
Gate~2 feasibility score.

\paragraph{MismatchAgent.}
Scores the severity of operator mismatch by comparing
$\mathbf{H}_{\text{nom}}$ against the expected
$\mathbf{H}_{\text{true}}$ using residual-norm heuristics and
registry-defined mismatch signatures.  Selects the appropriate
correction strategy from \texttt{solver\_registry.yaml}.

\paragraph{RecoverabilityAgent.}
Performs table-driven recoverability assessment: looks up the modality's
null-space characteristics, compression ratio, and known PSNR baselines.
Returns a Gate~1 pass/fail decision with a predicted PSNR range.

\paragraph{AnalysisAgent.}
Receives the three gate scores and performs bottleneck classification.
Identifies the binding gate per the Triad Law (\cref{sec:triad}) and
generates ranked suggestions for improving reconstruction quality.

\paragraph{Negotiator.}
Implements cross-agent veto logic.  Computes the joint probability of
success across all three gates and enforces a configurable threshold
below which the run is aborted with a diagnostic message rather than
producing a misleading reconstruction.

\paragraph{PreFlightReportBuilder.}
Assembles the final pre-flight report: expected PSNR, dominant gate,
correction strategy, estimated runtime, and resource requirements.  This
report is presented to the user (or downstream system) before committing
compute.

% ============================================================
%  Executable contract
% ============================================================
\subsection{ExperimentSpec: The Executable Contract}

The \texttt{ExperimentSpec}~v0.2.1, implemented as a Pydantic
\texttt{StrictBaseModel} with \texttt{extra="forbid"}, serves as the
executable contract between agents.  Every field is typed and validated;
NaN and Inf values are rejected at parse time.  When an LLM is used, it
returns \emph{only} registry IDs --- never raw parameters --- which are
then resolved and validated mechanically.

The contract ecosystem comprises \textbf{25~Pydantic models}, all
inheriting from \texttt{StrictBaseModel}, ensuring that no unvalidated
data enters the pipeline.  These models are backed by \textbf{9~YAML
registries} totaling \textbf{7,034~lines}, covering modalities, solvers,
noise profiles, mismatch types, correction strategies, metrics, hardware
profiles, datasets, and experiment templates.
